\documentclass[a4paper]{scrartcl}

\usepackage{csquotes}
\usepackage{amsmath}
\usepackage{amssymb}
\usepackage{amsthm}
\usepackage{graphicx}
\usepackage{float}
\usepackage{amsfonts}
\usepackage{subcaption}
\usepackage{url}
\usepackage{hyperref}
\usepackage{xurl}
\usepackage{cleveref}
\usepackage[labelfont=bf]{caption}
\usepackage{bbm}
\usepackage{mhchem}
\usepackage{units}

\bibliographystyle{vancouver}

% #1: path to the image
% #2: caption for the figure
% #3: label for the figure
% #4: optional height (default 7cm)
\newcommand{\insertfigure}[4][7cm]{%
    \begin{figure}[H]
        \centering
        \includegraphics[height=#1, width=\linewidth, keepaspectratio]{#2}
        \caption{#3}
        \label{#4}
    \end{figure}
}


% #1: path to the first image
% #2: path to the second image
% #3: caption for the entire figure
% #4: label for the figure
\newcommand{\insertfigures}[5][7cm]{
	\begin{figure}[H]
		\centering
		\begin{subfigure}[t]{0.49\textwidth}
			\centering
			\caption{}
			\includegraphics[height=#1, width=\linewidth,keepaspectratio]{#2}
		\end{subfigure}
		\hfill
		\begin{subfigure}[t]{0.49\textwidth}
			\centering
			\caption{}
			\includegraphics[height=#1, width=\linewidth, keepaspectratio]{#3}
		\end{subfigure}
		\caption{#4}
		\label{#5}
	\end{figure}
}

\pagenumbering{gobble}

\begin{document}

\begin{center}
    {
        Written report for \enquote{Machine Learning for Water Distribution Networks}\\
    }

    \vspace{1.5cm}

    {\large
        Report by\\
        Yannik Schroeder
    }

    \vspace{1.5cm}

    \today

    \vspace*{\fill}
\end{center}

\clearpage
\pagenumbering{Roman}

\tableofcontents

\pagebreak

\pagenumbering{arabic}

\section{Introduction}

\subsection{Motivation}
Reaction networks and the dynamics of particle concentrations over time play an important role in chemistry and chemical engineering~\cite{wikipedia_gnn}. Accurate physical simulations of large systems tend to be slow and resource-intensive. Graph neural networks (GNN) are a modern approach to solve the problem of predicting concentration changes in reaction networks. One of the biggest challenges when using GNNs is the generation of high quality training data. In this project, training data for reaction networks of 2D particles will be created using a custom C++ simulator called \textit{Cell} supporting simple uni- and bimolecular reactions. While not physically accurate, due to the simple design the simulation is relatively fast and supports several parameters similar to real-world systems. This makes it a suitable tool for quickly generating heterogenous datasets that can be used to train GNNs in order to understand how model topology and training parameters impact prediction accuracy, which will be the main focus of this project.

\subsection{The simulation}
\subsubsection{Configurations}
\textit{Cell}~\cite{cell} is a simple, collision-based reaction network simulator. The various properties of a simulation can be saved as a \path{.json} configuration file using the GUI. The GUI can run the simulation and display the current state and particle counts in real time. The main output of the simulation are the particle counts over time, as visible in the right-hand side of the screenshot in \cref{fig_cell_screenshot}.

\insertfigure{figures/cell_screenshot.png}{Screenshot of the cell simulation software. The simulation widget on the left shows the simulation state in real time. The plot widget on the right-hand side shows the disc counts over time.}{fig_cell_screenshot}

A CLI of the simulation is also available. It takes the configuration file path and the desired simulation duration as parameters and creates a \path{.csv} file containing the particle counts evaluated at each simulation time step. Since the configuration is the main input of the simulation, it also served as the main input of the deep learning model. It consists of the following parameters.

\subsubsection*{Discs}
The particles in the simulation are called \textit{discs}. The physical properties of discs are determined by their disc type $d = \{\text{name}, r, m\}$, where $r$ and $m$ are the radius (in px) and mass (dimensionless) of the type. To conserve total disc area, $r = \sqrt{m}$ was set for all disc types in the context of this project. The set of all disc types is called $D = \{d_1, \hdots, d_M\}$.

\subsubsection*{Simulation size and distribution}
The simulation runs in a circular confinement of radius $R$ (in px) with a total number of $N$ discs. Once disc types are defined, the initial disc type distribution $f = \{f_1, \hdots, f_M\}$ with $M = |D|$ can be specified. It assigns a relative frequency to each disc type. Using the largest disc type radius as the edge length, a square lattice spanning the confinement is calculated and discs according to their relative frequencies are randomly distributed across the grid until either the desired count $N$ or the total lattice capacity is reached.

\subsubsection*{Reactions}
The simulation supports four different reaction types: Transformation reactions $A\rightarrow B$, decomposition reactions $A \rightarrow B + C$, combination reactions $A+B \rightarrow C$ and exchange reactions $A+B \rightarrow C + D$. The simulation does not allow reactions where educt and product masses are not equal. Additionally, conservation of kinetic energy is guaranteed using velocity scaling where necessary. The unimolecular transformation and decomposition reactions are given a time based reaction probability $p$ per second:

\begin{align*}
    p = 1 - (1 - p')^{\frac{1}{\Delta t}} \Leftrightarrow p' = 1 - (1 - p)^{\Delta t},
\end{align*}

where $\Delta t$ is the reaction time step and $p'$ is the probability of the reaction to occur at each discrete reaction step. For example, if $p = 0.5$, $50\%$ of the educt will have reacted on average after one second. The bimolecular combination and exchange reactions have collision-based reaction probabilities $p$. For example, if $p = 0.5$, each collision of the educts has a $50\%$ probability to result in a reaction. If a set of educts has multiple reactions assigned to it, all reactions are tried sequentially in random order at each simulation time step or collision.

Each reaction can be represented as a set $\{\text{Educts}, \text{Products}, p\}$. The set of all reaction is called $\Omega$.

\subsubsection*{Most probable speed}
The velocity components of the discs are each sampled from a Gaussian distribution:

\begin{equation*}
    v_x, v_y \sim \mathcal{N}(0, \sigma^2).
\end{equation*}

The standard deviation is calculated using the user-defined most probable speed $v_{\text{mp}}$ and the mass $m$ of the disc:

\begin{equation*}
    \sigma = \frac{v_{\text{mp}}}{m}.
\end{equation*}

\subsubsection*{Time step}
The time step $\Delta t$ controls the accuracy of the simulation. After each simulation step, every disc is moved by $\Delta \vec{r} = \vec{v}\Delta t$. Smaller $\Delta t$ values make the simulation more accurate but are more CPU intensive. In this project, $\Delta t = 3\text{ms}$ for all simulations.

\subsubsection*{The configuration}
With the above notation, each simulation config can be written as a tuple

\begin{equation*}
    (D, R, N, f, \Omega, v_{\text{mp}}, \Delta t).
\end{equation*}

Since this project focusses on reaction networks, $R$, $N$, $v_\text{mp}$, $\Delta t$ and the simulation duration remained constant across all simulations to reduce the parameter space.

\subsection{Target reactions}
The target reaction network that drove dataset creation and training approaches is a toy network inspired by the formation of calcium carbonate in water. The network contains ten reactions, including all four different reaction types, and a total of nine different species. The graph of the network (see \cref{subsection_model_topology}) is shown in \cref{fig_water_network_reaction_graph}. The model created in this project is supposed to be a general-purpose model that can predict trajectories for any reaction network up the size of the one from \cref{fig_water_network_reaction_graph}.

\section{Related work}
Since the used simulator is a custom project, there are no related projects of exactly this kind. However, GNNs and their applications in chemical reaction networks have already been explored in various instances. For example, \cite{phygnet} describes a GNN which \enquote{maps microbial reaction networks into a directed graph by representing species and reactions as nodes and edges}~\cite{phygnet_website}, which is very similar to what was done in this more theoretical project.

\section{Methodology}
\subsection{Model topology}
\label{subsection_model_topology}
The main purpose of the model is to predict the trajectories of each species fraction over time. To do this, the model receives the disc types, their properties and current fractions, and the reactions as input. In this context, a fraction is the current number of discs of a given disc type, divided by the initial total disc count $N$. This normalization is useful to make trajectory plots comparable and to avoid large gradients during optimization, since absolute counts typically range from zero to several thousands.

The input of the model can be visualized as a graph. For example, a small graph containing all four reaction types is shown in \cref{fig_simple_network_graph}.

\insertfigure[10cm]{figures/simple_example_network.jpg}{A simple reaction network containing all four reaction types. The disc type nodes to the left (educts) and right (products) contain the disc type names and their masses as labels. The reaction nodes in the middle connect educts and products and contain the reaction itself and its probability as labels. The initial fractions are not shown in this figure.}{fig_simple_network_graph}

The current fractions for each species are also part of the graph, but vary after every time step. Thus, a single simulation run can be thought of as a tuple of $M$ graphs, where $M$ is the number of time steps simulated.

Reaction networks pose two distinct challenges for the model topology. First, the model should be able to deal with varying numbers of species and reactions. Second, permutations of the network such as reordering of the species or changing $A+B\rightarrow C$ into $B+A\rightarrow C$ should naturally be supported. A MLP where each input node corresponds to a simulation parameter, for example, would thus not be suitable. Since reaction networks can be naturally represented as graphs, a model suitable to operate on graphs as input was developed instead. Such models are called graph neural networks (GNNs).

The GNN used in this project predicts the flux from educt to product using two separate MLPs that operate on the level of reaction nodes: A reactant MLP $3 \rightarrow 32$ that calculates a message for each educt and a product MLP $35 \rightarrow 32 \rightarrow 1$ that actually predicts the flux from educt to product. All layers used the ReLU activation function. The final output was then fed into the sigmoid function to obtain a value in $[0, 1]$.

\subsection{Model prediction}
A prediction of the flux for a single reaction with a variable number of educts and products follows the following steps. Here, $f_{e_i}$ and $f_{p_i}$ are the educt and product fractions, $m_{e_i}$ and $m_{p_i}$ are the masses of the educt and product disc types and $p$ is the reaction probability.

\begin{enumerate}
    \item For each educt, calculate a message $\vec{M}$ using $(f_{e_i}, m_{e_i}, p)$ as input for the reactant MLP
    \item Calculate the sum $\vec{S}$ of all educt messages
    \item Use $(\vec{S}, f_{p_1}, f_{p_2}, p)$ as input for the product MLP. The MLP outputs a single number in $[0, 1]$ representing the predicted flux $r_p$
    \item Calculate the product of the educt fractions $P = f_{e_1}f_{e_2}$
    \item Calculate the requested flux $r$ using $r = r_p\cdot P \cdot p$
    \item Scale the requested flux such that it can't use up more educt than available
    \item Subtract the scaled requested flux from the educt fractions and add it to the product fractions
\end{enumerate}

Using the product $P$ of educt fractions was inspired by chemical kinetics, where reaction rates are usually proportional to the product of educt concentrations. The process of calculating a message for each educt and calculating a sum to pass to the next MLP is called \textit{message passing}~\cite{wikipedia_gnn}. Since addition is commutative, this solves the issue of order-dependent representations and supports an arbitrary number of educts. The same could have also been done for the products to avoid order-dependency of the product fractions.

\subsection{Model training}
When operating on a graph with multiple reactions, the above calculations for each step are done in parallel utilizing \path{torch.tensor} objects. During training, multiple graphs were batched together and indices were kept in sync with the help of the \path{data.Data} class from \path{torch_geometric}, a package specialized for working with GNNs~\cite{fey2019}. The loss was calculated using the mean squared error (MSE) of the prediction compared to the simulated trajectory.

In order to predict a trajectory, the model was repeatedly applied to the last predicted state, making it an autoregressive model. If the loss during training was calculated using only a single prediction, the model would likely struggle to learn that small initial errors can accumulate to large drifts over time. To this end, a \textit{rollout window} of one second was used during training: All simulated trajectories were resampled to $\Delta t = \unit[50]{ms}$ and the model was applied 20 times to the last predicted state. Each graph contained both the current species fractions and the fractions for each of the next 20 time steps. The total loss used for optimization was the sum of all 20 losses for each prediction. Training used the Adam optimizer~\cite{adam} with a learning rate of $1 \times 10^{-3}$, a batch size of $32$ and five epochs. Each epoch took about one to two hours to complete. The dataset for training was divided with a 90/10 split for training and testing.

\subsection{Datasets}
The limiting factor for dataset creation turned out to be RAM usage. A dataset containing about 2,000 CSV files of simulated trajectories already used about $\unit[50]{GB}$ of RAM even after resampling to $\Delta t = \unit[50]{ms}$. Since a linux server with about $\unit[8]{GB}$ of RAM and four $\unit[2.4]{GHz}$ processors was used for training, dataset creation using \path{torch_geometric} was done using a \unit[100]{GB} SWAP file and took about \unit[12]{h}.

To have the model learn dynamics of all types of reaction systems, the 2,000 graphs were divided evenly into four sections. Since there are four reaction types, each section contained an equal number of runs for all possible combinations without replacements for $n \in \{1, 2, 3, 4\}$ reaction types. For example, the first section contained 125 graphs for each reaction type, whereas the second contained about 83 graphs for each of the 6 possible combinations of two reaction types.

Each graph parameter was created by drawing from appropriate random distributions.

The number of disc types $M$ was drawn from $\mathcal{U}\{3, \hdots, 10\}$. $\mathcal{U}$ denotes the uniform distribution. With $M=3$, all four reactions are possible. Since reactions conserve mass, the masses were sampled by choosing a quantum $q \sim \mathcal{U}\{1, \hdots, \left\lceil \frac{100}{M - 1} \right\rceil\}$ and defining initial multipliers $\{1, 1, 2\}$. At least two identical masses and one twice as large guarantee that all reactions are possible. Then, $M-3$ new multipliers were added by randomly choosing a multiplier, and adding its value plus one to the list of multipliers. The masses were then calculated by multiplying $q$ with each multiplier.

To add reactions, the set of all possible reactions for each reaction type, taking mass conservation into account, was created first. To make each reaction type equally likely, each set was truncated to the size of the smallest set, and reactions were randomly added from each set using an constant edge probability. To get differently sparse graphs, the edge probability was randomly chosen from $\{\frac{2}{Z}, \frac{5}{Z}, \frac{10}{Z}\}$, where $Z$ is the total number of reactions after truncation. This creates graphs with either 2, 5 or 10 reactions on average by iterating all reactions and either skipping or adding them with the edge probability. Species not part of any reactions were removed from the graph.

Reaction probabilities in the target reaction network (\cref{fig_water_network_reaction_graph}) were mostly either close to zero or one. To reflect this in the training data, reaction probabilities were drawn from a randomly selected truncated normal distribution with parameters $(\mu,\sigma) \in \{(0.1,0.05),(0.5,0.2),(0.9,0.05)\}$, with respective selection probabilities $0.7$, $0.1$, and $0.2$. The resulting distribution is shown in \cref{fig_reaction_prob_distribution}.

\insertfigure{figures/reaction_prob_distribution.jpg}{Distribution of the reaction probabilities in the dataset obtained by sampling $10^6$ samples. The probability is sampled by choosing from one of three truncated normal distributions with different probabilities.}{fig_reaction_prob_distribution}

The initial fractions were randomly drawn from a Dirichlet distribution with $\alpha \sim \mathcal{U}[0.1, 3]$. Drawing $\alpha$ randomly assures that both graphs with about equal fractions for all disc types and graphs with a single species dominating are present in the dataset.

Two datasets were created using the same set of graphs: For the first, each simulation was run once, and for the second, the average trajectory of 10 runs was calculated.

\section{Results}
The target reaction network from \cref{fig_water_network_reaction_graph} has many disc types, but the formation of $\ce{CaCO3}$ is the most relevant aspect. Its predicted fraction and the fraction of some other educts alongside the simulation runs is shown in \cref{fig_water_prediction}.

\insertfigure{figures/water_prediction.jpg}{Prediction of $\ce{CaCO3}$ fractions over time using the model trained on single simulation runs to the left and using the model trained on averaged simulation runs to the right. The mean squared error (MSE) was calculated with the averaged simulation run as a comparison. The simulated run on the left was randomly chosen, the one on the right is the average of 10 runs}{fig_water_prediction}

Both models were able to predict the overall dynamic of the system up to a constant offset. For the simple network seen in \cref{fig_simple_network_graph}, the predictions alongside simulation runs are shown in \cref{fig_simple_prediction}. To test if input order affected the model output, various permutations of the reaction input order and reaction educt and product order were tested. The MSE differences between these predictions were $< 10^{-5}$.

\insertfigure{figures/simple_prediction.jpg}{Prediction of the simple network from \cref{fig_simple_network_graph} with the model trained on single simulation runs to the left and the model trained on averaged runs to the right. The MSE was calculated with the averaged runs as a reference. The left figure shows a single simulation run as comparison, the right the average of 10 runs.}{fig_simple_prediction}

For this simple network, the model failed to predict the dynamic of the trajectories.

The evaluation loss on the testing dataset of averaged runs was $1.47\cdot 10^{-4}$ for the model trained on single runs and $1.49\cdot 10^{-4}$ for the model trained on averaged runs. The model trained on averaged runs was also trained for 10 and 20 epochs as a comparison. The testing losses were $1.74\cdot 10^{-4}$ and $1.73\cdot 10^{-4}$ respectively.

\section{Discussion}
Both models trained on single and averaged runs performed roughly the same and were able to predict the trajectory of the relatively complex target reaction network from \cref{fig_water_network_reaction_graph}. Additionally, permutations of the input did not meaningfully impact model output, indicating that the examined prediction task is either not sensible to permutations or that the message passing for reaction educts successfully prevented order-dependent outputs. However, they were unable to predict correct trajectories for the much smaller network from \cref{fig_simple_network_graph}. This is most likely due to how sensible network dynamics are to small changes in parameters. For example, changing a single reaction probability in \cref{fig_water_network_reaction_graph} results in very different dynamics, as is shown in \cref{fig_simple_network_comp}. 

\insertfigure[10cm]{figures/simple_network_comp.jpg}{Averaged runs of the network from \cref{fig_simple_network_graph} with different probabilities for the decomposition reaction $C \rightarrow A + B$. The trajectories on the right correspond to 10 averaged runs with $p=0.3$, whereas on the left it was $p=0.15$. A small change in a single reaction can lead to completely different dynamics.}{fig_simple_network_comp}

\Cref{fig_simple_network_comp} indicates that the dataset should not only contain many different graphs, but also simulations of the same reaction network where only some parameters like reaction probabilities or masses were changed. Overall, many different parameters could contribute to the poor performance regarding individual reaction networks, not only the training dataset quality. For example, the requested educt flux was calculated by multiplying the product MLP model output with the reaction probability and the educt fractions. It is not immediately clear how altering this computation, for example by removing the reaction probability from the product, impacts model performance.

Interestingly, using averaged simulation runs instead of single runs for training did not lead to a significant improvement (see MSEs in \cref{fig_water_prediction,fig_simple_prediction}). The model is able to learn expected mean trajectories even if single training trajectories are mostly random. This indicates that stochastic noise of individual runs is not a limiting factor of the model performance. Increasing the number of epochs for training from the small initial number of five did also not lead to a significant improvement. This suggests that there are underlying limitations of the GNN architecture and the two MLPs are maybe not capable of representing the information of the datasets any better.

\section{Conclusion}
The target reaction network dynamic was, to a constant offset, correctly predicted. The models were, however, not capable of generalizing to arbitrary networks up to the size of the target reaction network. This is likely due to how small configuration changes can lead to different trajectories, which was insufficiently taken into account during dataset creation.

Future work could use an updated simulator with more realistic reaction physics, incorporating electrical forces and thermodynamics. Additionally, other simulation parameters like the absolute disc counts, simulation size and velocity distribution could be included in the training dataset. 

\pagebreak

\appendix

\section{Supplementary figures}

\setcounter{figure}{0}
\renewcommand\thefigure{\thesection.\arabic{figure}}

\insertfigure[17cm]{figures/water_network.jpg}{A more complicated reaction network inspired by the formation of calcium carbonate in water. Initial fractions were $\{f_{\ce{H2O}} = 0.8, f_{\ce{Ca^{2+}}} = 0.1, f_{\ce{CO2 (g)}} = 0.1\}$. Through a chain of reactions, $\ce{Ca^{2+}}$ will slowly be used up to form $\ce{CaCO3}$.}{fig_water_network_reaction_graph}

\pagebreak

\addcontentsline{toc}{section}{References}

\bibliography{refs}

\end{document}